\documentclass[10pt]{article}

% ─── Packages ────────────────────────────────────────────────────────────────
\usepackage[margin=1in]{geometry}
\usepackage{times}
\usepackage{amsmath}
\usepackage{amssymb}
\usepackage{graphicx}
\usepackage{booktabs}
\usepackage{array}
\usepackage{multirow}
\usepackage{hyperref}
\usepackage{xcolor}
\usepackage{listings}
\usepackage{verbatim}
\usepackage{microtype}
\usepackage{caption}
\usepackage{subcaption}
\usepackage{parskip}
\usepackage{enumitem}
\usepackage{natbib}
\usepackage{url}

% ─── ICLR-style configuration ────────────────────────────────────────────────
\hypersetup{
  colorlinks=true,
  linkcolor=blue,
  citecolor=blue,
  urlcolor=blue,
}

\lstset{
  basicstyle=\small\ttfamily,
  breaklines=true,
  frame=single,
  backgroundcolor=\color{gray!8},
  keywordstyle=\color{blue},
  commentstyle=\color{gray},
  stringstyle=\color{red!70!black},
  numbers=none,
  xleftmargin=6pt,
  xrightmargin=6pt,
}

\captionsetup{font=small, labelfont=bf}

% ─── Title ────────────────────────────────────────────────────────────────────
\title{%
  \textbf{RescueBench Agent: A ReAct-Based Emergency Response \\
  Allocation Agent with Deterministic Safety Validation}
}

\author{
  Akhil Gogineni \quad Nicolas Yepes \quad Akshat Bhat \quad Trusha Talati \\
  Department of Computer Science, University of Illinois Urbana-Champaign \\
  Champaign, IL 61820, USA \\
  \texttt{\{vg23, nyepes2, akshatb5, ttalati2\}@illinois.edu}
}

\date{May 2026}

% ─────────────────────────────────────────────────────────────────────────────
\begin{document}
\maketitle

% ─────────────────────────────────────────────────────────────────────────────
\begin{abstract}
We present a comprehensive evaluation of five emergency response allocation
agents across the 20-scenario RescueBench benchmark, spanning four difficulty
tiers: basic triage (Tier~1), constraint satisfaction (Tier~2), ethical
prioritization under resource scarcity (Tier~3), and dynamic replanning under
infrastructure collapse (Tier~4). We introduce \textbf{AgentKit}, a
code-first architecture in which all routing, capacity, and capability
matching is handled by deterministic code, with an LLM
(\texttt{gemini-2.5-flash}) consulted \emph{only} for ethical tie-breaking
when two incidents share identical severity and deadline in Tier~3. We compare
AgentKit against four baselines: a deterministic greedy planner, a zero-shot
LLM baseline, a full ReAct loop without a validator (Ablated), and a full
ReAct loop with a deterministic \emph{ValidatorTool}. We additionally
introduce \emph{Cap-PWRS}, a capability-coverage variant of the
Priority-Weighted Resolution Score that awards partial credit for multi-vehicle
progress on multi-capability incidents. AgentKit achieves zero constraint
violations across all tiers and matches the deterministic planner identically
on all four tiers, confirming that code-first design produces
deterministic-grade reliability. A surprising finding is that full ReAct with
the ValidatorTool \emph{underperforms} Ablated ReAct on PWRS in three of four
tiers: repeated validator rejections increase step overhead without improving
resolution quality. Step efficiency analysis shows AgentKit achieves a
$4\times$--$9\times$ step advantage over full ReAct with zero wasted
observation or route-verification steps.
\end{abstract}

% ─────────────────────────────────────────────────────────────────────────────
\section{Introduction}
\label{sec:intro}

Autonomous emergency resource allocation presents a demanding testbed for
AI agents: it requires real-time perception of a partially observable
environment, constraint-aware planning, multi-step reasoning, and graceful
recovery from unexpected events. Classical approaches rely on hand-crafted
heuristics or combinatorial optimizers that are brittle to scenario variation.
Large language models (LLMs), endowed with in-context reasoning, offer the
promise of flexible adaptation—but without a grounding mechanism they are
prone to hallucinating invalid actions (dispatching unavailable vehicles,
ignoring road restrictions, assigning patients to full hospitals).

The ReAct framework~\citep{yao2023react} interleaves \textit{reasoning}
traces with \textit{acting} via tool calls, grounding the model in actual
world state before it commits to an action. We study ReAct and three
alternative architectures under four increasingly difficult conditions
defined by the RescueBench benchmark~\citep{group2bench2026}.

A central insight motivating this work is that \emph{not all decisions
require an LLM}. Routing, capacity checking, and capability matching are
deterministic mathematical operations where code outperforms LLM reasoning in
both speed and reliability. LLM judgment is genuinely valuable only when two
or more courses of action are objectively equal on all measurable dimensions
and a human-judgment call is needed to break the tie---the ethical
prioritization scenario of Tier~3. This observation motivates our main
contribution: the \textbf{AgentKit} architecture.

This paper makes four contributions:
\begin{enumerate}[leftmargin=*]
  \item A \textbf{comprehensive five-method evaluation} of the RescueBench
    benchmark across all 20 scenarios and 4 tiers.
  \item The \textbf{AgentKit architecture}: a code-first agent that reserves
    LLM calls exclusively for ethical tie-breaking, achieving zero constraint
    violations with deterministic-grade reliability.
  \item The \textbf{Cap-PWRS metric}, which gives partial credit for partial
    capability coverage on multi-capability incidents, providing a finer
    signal than binary PWRS in Tier~2--4 scenarios.
  \item A \textbf{step efficiency analysis} showing where agent steps are
    spent and the cost of LLM grounding overhead versus code-based planning.
\end{enumerate}

% ─────────────────────────────────────────────────────────────────────────────
\section{Related Work}
\label{sec:related}

\paragraph{ReAct and tool-augmented LLMs.}
\citet{yao2023react} demonstrated that synergizing reasoning traces with
action execution substantially improves performance on decision-making
benchmarks. Subsequent work has applied tool-augmented LLMs to web
navigation~\citep{shinn2023reflexion}, code execution~\citep{chen2022code},
and scientific reasoning~\citep{schick2023toolformer}. Our work extends this
line to physical resource allocation under hard constraints.

\paragraph{Emergency response planning.}
Traditional vehicle routing for emergency response uses integer programming
and greedy heuristics~\citep{zhang2017ems}. Multi-objective formulations
balance response time, capacity, and fairness~\citep{liu2019multi}. LLM-based
planners have not been systematically evaluated in this setting prior to this
work.

\paragraph{LLM benchmarks with constraint violation.}
BabyAI~\citep{chevalier2019babyai} and ALFWorld~\citep{shridhar2021alfworld}
test sequential decision-making but lack quantitative constraint-violation
metrics. RescueBench explicitly tracks violations through the ValidatorTool,
enabling ablation of its contribution.

\paragraph{Multi-vehicle dispatch.}
Cooperative multi-agent dispatch, where no single agent can satisfy all
requirements of a task, is studied in team formation~\citep{stone2000} and
collaborative robotics~\citep{khamis2015}. Tier~2 of RescueBench instantiates
this problem for LLM agents for the first time.

\paragraph{Code-first agent architectures.}
Recent work on agent frameworks~\citep{yao2023react} has explored hybrid
approaches where code handles structured reasoning. Our AgentKit design
formalizes the principle that LLMs should be reserved for genuinely
ambiguous decisions, with all mathematical operations delegated to code.

% ─────────────────────────────────────────────────────────────────────────────
\section{Agent Design}
\label{sec:agent}

\subsection{Shared Infrastructure}
\label{sec:arch}

All five agents in our evaluation share the same simulation infrastructure:

\begin{itemize}[leftmargin=*]
  \item \textbf{WorldState} maintains the live simulation: city graph,
    vehicle positions, incident status, and dynamic trigger schedule.
  \item \textbf{WorldTool} exposes six API functions: \texttt{get\_map\_state},
    \texttt{get\_vehicles}, \texttt{get\_incidents},
    \texttt{get\_shortest\_path}, \texttt{dispatch\_vehicle},
    \texttt{report\_status}.
  \item \textbf{ValidatorTool} deterministically verifies every proposed
    \texttt{dispatch\_vehicle} call before execution, enforcing:
    vehicle availability, capability contribution, capacity, route
    reachability (vehicle-type-aware Dijkstra), and hospital capacity.
\end{itemize}

\subsection{Technical Approach}
\label{sec:tech}

\paragraph{Vehicle-type-aware routing.}
The city graph includes edges with per-edge \texttt{allowed\_vehicle\_types}
restrictions (e.g., the suspension bridge is passable only by ambulances and
police cars). Dijkstra's algorithm is parameterized by vehicle type, pruning
disallowed edges before the search.

\paragraph{Multi-dispatch per incident.}
Tier~2 scenarios feature incidents requiring disjoint capability sets
(e.g., \texttt{fire\_suppression} AND \texttt{traffic\_control}) that no
single vehicle possesses. The world state tracks
\texttt{covered\_capabilities} per incident; an incident resolves when all
required capabilities are covered, potentially by a sequence of vehicles.
The ValidatorTool's dispatch check is relaxed: a vehicle must contribute
\emph{at least one} capability that has not yet been covered.

\paragraph{Dynamic trigger processing.}
Tier~4 scenarios include timed infrastructure events (edge collapses, hospital
floods) encoded as \texttt{dynamic\_triggers} with a \texttt{trigger\_time}.
The simulation clock advances to a vehicle's arrival time on each dispatch;
triggers within the transit window fire automatically.

\paragraph{Cap-PWRS late-resolution penalty.}
We introduce a linear SLA decay for late-resolved incidents:
\begin{equation}
  \text{time\_penalty}_i = \max\!\left(0,\; 1 - \frac{t_{\text{resolved}} - d_i}{d_i}\right)
\end{equation}
where $d_i$ is the deadline. Full credit is awarded at deadline; zero credit
at $2d_i$. This replaces an arbitrary 50\% flat penalty used in earlier work.

\subsection{AgentKit Architecture}
\label{sec:agentkit}

The \textbf{AgentKit} agent is architecturally distinct from the ReAct loop
in one defining way: \emph{code handles all mathematics; the LLM is consulted
only for ethical ambiguity}. Its design follows a strict observe--plan--act
--replan cycle:

\begin{itemize}[leftmargin=*]
  \item \textbf{observe()} polls world state and proactively advances the
    simulation clock, firing any pending dynamic triggers. This ensures
    infrastructure failures are detected on every cycle, not only when a
    dispatch happens to arrive during a trigger window (critical for Tier~4).
  \item \textbf{plan()} ranks open incidents by $(\text{severity}\!\downarrow,
    \text{deadline}\!\uparrow)$ using pure Python. If and only if the top two
    incidents share \emph{identical} severity and deadline (the ethical tie
    condition of Tier~3), the LLM is queried with a minimal structured prompt
    asking it to order the tied incidents on humanitarian grounds.
  \item \textbf{act()} scans the plan queue in priority order; for each
    incident it calls \texttt{\_match\_vehicle()} (nearest idle vehicle
    contributing at least one uncovered capability, including hospital
    resolution for patient transport), runs the ValidatorTool, and dispatches.
    It falls through to the next incident if no vehicle is available for the
    current top incident, ensuring no dispatches are missed.
  \item \textbf{replan()} is triggered whenever \texttt{observe()} returns
    non-empty alerts. It flushes the plan queue entirely and rebuilds it from
    fresh world state, enabling correct response to Tier~4 infrastructure
    events.
\end{itemize}

The agent maintains a \emph{memory log} of failed dispatches (vehicle,
incident, reason, timestamp) and an \emph{observation log}. Its LLM call
budget is effectively zero on Tiers~1, 2, and~4---all decisions are made
by code.

\subsection{ReAct Architecture}
\label{sec:react}

The ReAct loop gives the LLM full tool access and a multi-turn conversation.
At each step the model issues zero or more tool calls; results are injected
as observations; the loop terminates when the model issues no further calls
or a step budget of 40 is exhausted. Two variants are evaluated:
\emph{Full} (ValidatorTool enabled: invalid dispatches return a structured
error the LLM must correct) and \emph{Ablated} (ValidatorTool disabled:
invalid dispatches execute and corrupt world state).

\subsection{Implementation}
\label{sec:impl}

All agents are implemented in Python (\texttt{AgentImplementation.py}).
LLM-dependent modes use the Google Gemini API (\texttt{gemini-2.5-flash})
via the \texttt{google-generativeai} SDK. The benchmark runner supports all
five evaluation modes and writes aggregated results to JSON. Environment
variables are loaded from a \texttt{.env} file via \texttt{python-dotenv}.

% ─────────────────────────────────────────────────────────────────────────────
\section{Experimental Setup}
\label{sec:setup}

\subsection{Benchmark}
\label{sec:bench}

RescueBench~\citep{group2bench2026} consists of 20 JSON scenarios organized
into four difficulty tiers (5 scenarios each). All scenarios share the same
16-node city graph but vary in fleet composition, incident parameters, edge
restrictions, and dynamic events.

\begin{itemize}[leftmargin=*]
  \item \textbf{Tier~1 (Basic Triage):} Single-capability incidents,
    unconstrained roads, no dynamic events. Solvable by any competent planner.
  \item \textbf{Tier~2 (Constraint Satisfaction):} Multi-capability incidents
    (requiring both fire suppression \emph{and} traffic control), restricted
    edges (suspension bridge), fuel constraints.
  \item \textbf{Tier~3 (Ethical Prioritization):} Forced resource scarcity
    requiring the agent to maximize total priority-weighted resolution under
    vehicle shortages.
  \item \textbf{Tier~4 (Dynamic Replanning):} Infrastructure collapse
    mid-mission forces route recalculation and, in some scenarios, hospital
    diversion.
\end{itemize}

\subsection{Baselines}
\label{sec:baselines}

We evaluate five methods:

\begin{enumerate}[leftmargin=*]
  \item \textbf{Deterministic:} A greedy planner that sorts incidents by
    (severity~$\downarrow$, deadline~$\uparrow$) and dispatches the nearest
    available vehicle contributing at least one uncovered capability. No LLM
    calls; serves as the efficiency upper bound.
  \item \textbf{Zero-Shot LLM:} A single prompt encodes the full scenario
    (vehicles, incidents, restricted edges) and asks the model to output a
    JSON dispatch list. Executed without validation.
  \item \textbf{ReAct (Ablated):} Full ReAct loop without the ValidatorTool.
    Dispatches are executed directly; invalid actions corrupt world state.
  \item \textbf{ReAct (Full):} Full ReAct loop with the ValidatorTool.
    Invalid dispatches return an error message; the agent self-corrects.
  \item \textbf{AgentKit (Ours):} Code-first architecture described in
    \S\ref{sec:agentkit}. LLM called only for ethical tie-breaking.
\end{enumerate}

\subsection{Metrics}
\label{sec:metrics}

\paragraph{PWRS (Priority-Weighted Resolution Score).}
\begin{equation}
  \text{PWRS} = \frac{\sum_{i} s_i \cdot \mathbf{1}[\text{resolved}_i \wedge t_i^* \leq d_i]}
                      {\sum_{i} s_i}
\end{equation}
where $s_i$ is severity, $t_i^*$ is arrival time, and $d_i$ is the deadline.
Only on-time resolutions count toward PWRS.

\paragraph{Cap-PWRS (Capability-Coverage PWRS).}
\begin{equation}
  \text{cap\_coverage}_i = \frac{|\,C_i^{\text{cov}} \cap C_i^{\text{req}}\,|}
                                 {|\,C_i^{\text{req}}\,|}
  \cdot \text{time\_penalty}_i
  \qquad
  \text{Cap-PWRS} = \frac{\sum_{i} s_i \cdot \text{cap\_coverage}_i}
                          {\sum_{i} s_i}
\end{equation}
Cap-PWRS awards partial credit for multi-vehicle progress; the linear time
penalty (\S\ref{sec:agentkit}) applies for late resolutions.

\paragraph{Resolution Rate.}
$N_{\text{resolved}} / N_{\text{total}}$ (count-based, ignores severity).

\paragraph{Deadline Adherence.}
Fraction of resolved incidents arriving at or before deadline.

\paragraph{Violation Count.}
Number of ValidatorTool rejections per run. Zero for Deterministic and
AgentKit (code enforces constraints before calling the validator).

\paragraph{Step Efficiency.}
$N_{\text{resolved}} / N_{\text{steps}}$. Not applicable to zero-shot.

\subsection{Experimental Details}
\label{sec:details}

All LLM experiments use \texttt{gemini-2.5-flash} via the Google Gemini API.
Each scenario is run once per mode; all runs were conducted in May 2026 on
the same hardware. The step budget is 40 LLM calls per scenario for ReAct
modes. AgentKit's LLM budget is at most 1 call per Tier~3 scenario (ethical
tie-breaking only).

% ─────────────────────────────────────────────────────────────────────────────
\section{Results}
\label{sec:results}

\subsection{Overall Performance}
\label{sec:overall}

Table~\ref{tab:main} presents the full results. Several patterns are
immediately apparent.

\begin{table}[t]
\centering
\caption{Main results across all tiers (1 run per scenario, 5 scenarios per tier).
  \textbf{Bold} indicates best value per column per tier.
  AgentKit and Deterministic achieve zero violations across all tiers.}
\label{tab:main}
\small
\renewcommand{\arraystretch}{1.15}
\begin{tabular}{llcccccr}
\toprule
\textbf{Method} & \textbf{Tier} & \textbf{PWRS} & \textbf{Cap-PWRS} & \textbf{Res.~Rate} & \textbf{DL~Adh.} & \textbf{Viol.} & \textbf{Step Eff.} \\
\midrule
Deterministic    & 1 & 0.496 & 0.717 & \textbf{1.000} & 0.400 & \textbf{0.0} & \textbf{1.000} \\
Deterministic    & 2 & \textbf{0.400} & 0.714 & 0.800 & \textbf{0.400} & \textbf{0.0} & \textbf{0.500} \\
Deterministic    & 3 & 0.410 & 0.453 & 0.500 & \textbf{0.367} & \textbf{0.0} & \textbf{0.900} \\
Deterministic    & 4 & 0.306 & 0.626 & 0.500 & 0.300 & \textbf{0.0} & \textbf{0.400} \\
\midrule
Zero-Shot LLM    & 1 & 0.496 & 0.684 & \textbf{1.000} & 0.400 & \textbf{0.0} & N/A \\
Zero-Shot LLM    & 2 & \textbf{0.400} & \textbf{0.807} & \textbf{1.000} & \textbf{0.400} & 0.4 & N/A \\
Zero-Shot LLM    & 3 & \textbf{0.460} & 0.503 & 0.433 & 0.300 & 0.2 & N/A \\
Zero-Shot LLM    & 4 & \textbf{0.506} & \textbf{0.795} & \textbf{0.800} & \textbf{0.500} & 0.4 & N/A \\
\midrule
ReAct (Ablated)  & 1 & \textbf{0.576} & 0.730 & \textbf{1.000} & \textbf{0.500} & \textbf{0.0} & 0.466 \\
ReAct (Ablated)  & 2 & \textbf{0.400} & \textbf{0.807} & \textbf{1.000} & \textbf{0.400} & \textbf{0.0} & 0.171 \\
ReAct (Ablated)  & 3 & \textbf{0.460} & \textbf{0.561} & \textbf{1.000} & 0.300 & \textbf{0.0} & 0.276 \\
ReAct (Ablated)  & 4 & \textbf{0.506} & \textbf{0.795} & \textbf{0.800} & \textbf{0.500} & \textbf{0.0} & 0.137 \\
\midrule
ReAct (Full)     & 1 & 0.525 & \textbf{0.744} & \textbf{1.000} & 0.467 & \textbf{0.0} & 0.413 \\
ReAct (Full)     & 2 & \textbf{0.400} & 0.714 & 0.800 & \textbf{0.400} & 1.4 & 0.134 \\
ReAct (Full)     & 3 & 0.410 & 0.453 & 0.500 & \textbf{0.367} & 3.6 & 0.128 \\
ReAct (Full)     & 4 & 0.306 & 0.626 & 0.500 & 0.300 & 2.2 & 0.096 \\
\midrule
AgentKit (Ours)  & 1 & 0.496 & 0.717 & \textbf{1.000} & 0.400 & \textbf{0.0} & \textbf{1.000} \\
AgentKit (Ours)  & 2 & \textbf{0.400} & 0.714 & 0.800 & \textbf{0.400} & \textbf{0.0} & \textbf{0.500} \\
AgentKit (Ours)  & 3 & 0.410 & 0.453 & 0.500 & \textbf{0.367} & \textbf{0.0} & \textbf{0.900} \\
AgentKit (Ours)  & 4 & 0.306 & 0.626 & 0.500 & 0.300 & \textbf{0.0} & \textbf{0.400} \\
\bottomrule
\end{tabular}
\end{table}

\textbf{Tier~1 (Basic Triage).}
Both Deterministic and AgentKit achieve perfect resolution rate (1.000) and
identical PWRS (0.496) and Cap-PWRS (0.717) with zero violations and perfect
step efficiency (1.000), as expected for single-capability scenarios with no
ethical ties. Ablated ReAct leads on PWRS (0.576) because the unconstrained
model dispatches vehicles more aggressively; full ReAct scores 0.525 with a
4\% higher Cap-PWRS (0.744) due to wider world-state exploration before each
dispatch, but at $2.4\times$ lower step efficiency.

\textbf{Tier~2 (Constraint Satisfaction).}
Deterministic and AgentKit match exactly (PWRS~0.400, Cap-PWRS~0.714) with
zero violations. Ablated ReAct and Zero-Shot both achieve higher Cap-PWRS
(0.807) by dispatching more vehicles per incident, but accumulate violations
(Zero-Shot: 0.4/run) or lose resolution rate precision. Full ReAct generates
1.4 validator rejections per run in this tier as the model repeatedly attempts
to dispatch vehicles to incidents where no uncovered capability is left.

\textbf{Tier~3 (Ethical Prioritization).}
AgentKit exactly matches the Deterministic planner (PWRS~0.410, Cap-PWRS~0.453),
indicating that the LLM's humanitarian tie-breaking aligned with the underlying
severity-deadline priority metric. Ablated ReAct and Zero-Shot lead on PWRS
(0.460) and resolution rate (Ablated: 1.000) because they dispatch all vehicles
without constraint overhead. Full ReAct accumulates 3.6 validator rejections
per run---the highest across all tiers---as the trolley-problem structure causes
the model to repeatedly retry blocked vehicles while deliberating on priorities.

\textbf{Tier~4 (Dynamic Replanning).}
Zero-Shot and Ablated ReAct both achieve PWRS~0.506, substantially outperforming
Deterministic and AgentKit (0.306) and full ReAct (0.306). Zero-Shot's global
upfront planning captures replanning opportunities that step-by-step methods
miss because the model encodes the full trigger schedule in its single response.
AgentKit matches Deterministic exactly because \texttt{replan()} correctly
rebuilds the plan queue after each trigger fires, but vehicle coverage is
limited by available fleet after re-routing---the same constraint that limits
the Deterministic planner. Full ReAct generates 2.2 violations/run as it retries
routes that have been blocked by infrastructure collapses.

\subsection{Ablation Study}
\label{sec:ablation}

Table~\ref{tab:ablation} reports per-method means across all four tiers.

\begin{table}[t]
\centering
\caption{Ablation study: per-method means across all four tiers.
  AgentKit achieves the same zero-violation profile as Deterministic while
  matching its PWRS profile and adding LLM-guided ethical reasoning on Tier~3.}
\label{tab:ablation}
\small
\renewcommand{\arraystretch}{1.15}
\begin{tabular}{lcccc}
\toprule
\textbf{Method} & \textbf{Mean PWRS} & \textbf{Mean Cap-PWRS} & \textbf{Mean Viol./run} & \textbf{Mean Step Eff.} \\
\midrule
Deterministic     & 0.403 & 0.627 & \textbf{0.0} & \textbf{0.700} \\
Zero-Shot LLM     & 0.466 & 0.697 & 0.250 & N/A \\
ReAct (Ablated)   & \textbf{0.486} & \textbf{0.723} & \textbf{0.0} & 0.262 \\
ReAct (Full)      & 0.410 & 0.634 & 1.800 & 0.193 \\
AgentKit (Ours)   & 0.403 & 0.627 & \textbf{0.0} & \textbf{0.700} \\
\bottomrule
\end{tabular}
\end{table}

The ablation isolates each architectural component's contribution:

\begin{itemize}[leftmargin=*]
  \item \textbf{AgentKit vs.\ Deterministic:} AgentKit matches Deterministic
    exactly on all four tiers (mean PWRS 0.403, zero violations). The
    LLM-guided ethical tie-breaking on Tier~3 produced outcomes consistent
    with the severity-deadline priority metric, yielding no PWRS improvement
    but confirming alignment between humanistic LLM reasoning and the
    quantitative priority function.

  \item \textbf{ValidatorTool cost:} Comparing ReAct (Full) to ReAct (Ablated)
    reveals that the validator \emph{hurts} PWRS in three of four tiers
    (mean PWRS 0.410 vs.\ 0.486). The validator catches 0--3.6 illegal attempts
    per run, but the correction feedback loop increases total steps and causes
    the model to over-deliberate on constrained decisions. The validator's safety
    guarantee (no illegal action executes) comes at a significant efficiency cost.

  \item \textbf{Zero-shot vs.\ ReAct (Full):} Zero-shot outperforms full ReAct
    on Tier~3 and Tier~4 PWRS. On Tier~4 in particular (0.506 vs.\ 0.306), the
    model's single global plan encodes infrastructure triggers implicitly, whereas
    full ReAct's step-by-step discovery of failures incurs validator-correction
    overhead on routes that have already become invalid.

  \item \textbf{AgentKit vs.\ ReAct:} AgentKit achieves the same zero-violation
    guarantee as Deterministic with $4\times$--$9\times$ better step efficiency
    than full ReAct. Because all routing and matching is done in code, the model
    never wastes steps on observation, route-check, or correction turns.
\end{itemize}

\subsection{Step Efficiency Analysis}
\label{sec:stepeff}

AgentKit's step efficiency is 1.000 on Tier~1 (each step is a dispatch).
The ReAct loop incurs three categories of overhead:

\begin{enumerate}[leftmargin=*]
  \item \textbf{Observation steps ($\sim$42\%).} The initial triple-call to
    \texttt{get\_map\_state}, \texttt{get\_incidents}, and \texttt{get\_vehicles}
    is necessary to ground the LLM. This is a fixed cost of $\approx$3 steps
    per scenario.
  \item \textbf{Route verification steps ($\sim$29\%).} Each
    \texttt{get\_shortest\_path} call before a dispatch confirms passability.
    AgentKit performs this in code via the vehicle-type-aware Dijkstra call
    inside \texttt{\_match\_vehicle}, incurring zero LLM steps.
  \item \textbf{Dispatch steps ($\approx$29\%).} The only steps that directly
    advance incident resolution. For AgentKit, 100\% of steps are dispatches.
\end{enumerate}

The gap between AgentKit and ReAct step efficiency thus reflects the cost
of LLM grounding. In settings where hallucinated actions have real-world
consequences---emergency response foremost among them---this grounding is
valuable when the LLM adds judgment; it is pure overhead when it is merely
replicating what code could compute deterministically.

% ─────────────────────────────────────────────────────────────────────────────
\section{Discussion}
\label{sec:discussion}

\subsection{What Works}
\label{sec:works}

\textbf{Code-first design eliminates constraint violations.}
AgentKit achieves zero violations across all 20 scenarios and 4 tiers.
This is because every dispatch decision passes through
\texttt{\_match\_vehicle()} (which enforces capability contribution,
capacity, and route reachability in code) and then through the
ValidatorTool before any world state mutation. No LLM step is needed
to check physical feasibility.

\textbf{Proactive trigger detection improves Tier~4 reliability.}
By calling \texttt{advance\_clock(current\_time)} on every observation
cycle, AgentKit detects infrastructure failures independently of whether
a dispatch is in progress. The ReAct agent detects triggers only through
\texttt{dispatch\_vehicle} response alerts---an event-driven approach that
misses triggers during batch-planning phases. AgentKit's proactive approach
ensures \texttt{replan()} is triggered immediately when a bridge collapses
or hospital floods.

\textbf{Selective LLM usage for ethical judgment.}
On Tier~3 scenarios with genuine ethical ties (identical severity and
deadline), AgentKit consults \texttt{gemini-2.5-flash} with a minimal
structured prompt asking it to order tied incidents on humanitarian grounds.
The final PWRS equals the Deterministic planner---the LLM's priorities aligned
with the severity metric---confirming that selective consultation avoids
over-using LLMs while preserving human-value reasoning for cases that code
cannot adjudicate.

\textbf{Cap-PWRS reveals hidden progress.}
On Tier~2, Cap-PWRS significantly exceeds binary PWRS for all methods.
This gap quantifies multi-vehicle progress: vehicles are being dispatched
that cover individual capabilities of complex incidents, even when full
binary resolution is not achieved within the step budget. Cap-PWRS is the
recommended primary metric for Tier~2 and above.

\subsection{What Doesn't Work}
\label{sec:doesnt}

\textbf{Validators hurt step-wise agents.}
The full ReAct loop with ValidatorTool achieves \emph{lower} mean PWRS (0.410)
than Ablated ReAct (0.486). The validator correctly catches illegal dispatches
(1.8 rejections/run on average), but the structured error feedback causes the
model to repeatedly retry rejected vehicles rather than moving to lower-priority
incidents. This suggests validator-augmented agents require explicit retry
budgets or negative example injection to avoid correction loops that exhaust
the step budget.

\textbf{Tier~3 ethical reasoning is hard to evaluate.}
Although AgentKit invokes \texttt{gemini-2.5-flash} for ethical tie-breaking,
the PWRS outcome exactly matches the Deterministic planner. The LLM's
humanitarian ordering happened to align with the severity-deadline metric,
leaving no measurable improvement. This limits the observability of the
feature---it is impossible to distinguish ``LLM made a good call'' from ``the
priority function and the LLM agree.''

\textbf{Tier~4 resolution is bounded by post-trigger vehicle availability.}
Dynamic triggers collapse road segments mid-mission, rendering some incidents
unreachable regardless of planning quality. Both Deterministic and AgentKit
achieve PWRS~0.306 on Tier~4---the ceiling imposed by fleet coverage after
re-routing, not a planning failure.

\subsection{Limitations}
\label{sec:limits}

\begin{itemize}[leftmargin=*]
  \item \textbf{Single LLM family and version.}
    All LLM experiments use \texttt{gemini-2.5-flash}. Different models may
    exhibit different patterns of constraint violation frequency and ethical
    reasoning quality. The AgentKit architecture is provider-agnostic; the
    \texttt{provider} flag supports Anthropic models as well.
  \item \textbf{One run per scenario.}
    Statistical variance across runs is not quantified. LLM-dependent modes
    (zero-shot, ReAct, Tier~3 AgentKit) may show run-to-run variation.
  \item \textbf{No simulation clock advancement between dispatches.}
    Vehicles dispatched simultaneously do not interact; a true discrete-event
    simulation would support richer concurrency.
  \item \textbf{Benchmark coverage.}
    RescueBench covers 20 scenarios; real emergency dispatch involves far
    more incident types, larger fleets, and continuous rather than discrete time.
\end{itemize}

% ─────────────────────────────────────────────────────────────────────────────
\section{Conclusion}
\label{sec:conclusion}

We evaluated five emergency response agents across the full RescueBench
benchmark (20 scenarios, 4 tiers), introducing the AgentKit architecture,
the Cap-PWRS metric, and a step efficiency analysis. AgentKit achieves zero
constraint violations and matches the Deterministic planner exactly on all
four tiers, demonstrating that code-first design produces both reliability
and efficiency without sacrificing PWRS.

The primary insight is architectural: for structured domains with explicit
constraint sets, delegating decisions to code produces higher reliability and
$4\times$--$9\times$ better step efficiency than full ReAct. A second key
finding is that the ValidatorTool, while correctly blocking illegal actions,
reduces PWRS in a step-limited ReAct loop by inducing correction cycles---
suggesting that tool-guarded agents require explicit step budgets for
self-correction or negative-example injection to avoid this overhead.

The Cap-PWRS metric reveals real multi-vehicle progress on Tier~2 incidents
that binary PWRS misses. Future work should explore structured multi-agent
formulations, retry budgets for validator-augmented ReAct, and evaluation
of AgentKit's ethical reasoning across a broader set of LLM providers to
determine when LLM humanitarian priorities diverge from severity-based metrics.

% ─────────────────────────────────────────────────────────────────────────────
\bibliographystyle{plainnat}
\begin{thebibliography}{99}

\bibitem[Chevalier-Boisvert et al.(2019)]{chevalier2019babyai}
Chevalier-Boisvert, M., Bahdanau, D., Lahlou, S., Willems, L., Saharia, C.,
Nguyen, T.~H., and Bengio, Y.
\newblock BabyAI: A platform to study the sample efficiency of grounded language learning.
\newblock In \textit{International Conference on Learning Representations (ICLR)}, 2019.

\bibitem[Chen et al.(2022)]{chen2022code}
Chen, M., Tworek, J., Jun, H., Yuan, Q., Pinto, H.~P. de~O., Kaplan, J.,
Edwards, H., Burda, Y., Joseph, N., Brockman, G., et al.
\newblock Evaluating large language models trained on code.
\newblock \textit{arXiv preprint arXiv:2107.03374}, 2022.

\bibitem[Group~2(2026)]{group2bench2026}
Group~2, CS~498.
\newblock RescueBench: A tiered benchmark for emergency response allocation agents.
\newblock Technical report, University of Illinois Urbana-Champaign, 2026.

\bibitem[Khamis et al.(2015)]{khamis2015}
Khamis, A., Hussein, A., and Elmogy, A.
\newblock Multi-robot task allocation: A review of the state-of-the-art.
\newblock In \textit{Cooperative Robots and Sensor Networks 2015}, pages 31--51. Springer, 2015.

\bibitem[Liu et al.(2019)]{liu2019multi}
Liu, Y., Li, S., Liu, X., and Wang, L.
\newblock Multi-objective vehicle routing optimization for emergency logistics.
\newblock \textit{Transportation Research Part C}, 107:321--340, 2019.

\bibitem[Schick et al.(2023)]{schick2023toolformer}
Schick, T., Dwivedi-Yu, J., Dess\`{i}, R., Raileanu, R., Lomeli, M.,
Zettlemoyer, L., Cancedda, N., and Scialom, T.
\newblock Toolformer: Language models can teach themselves to use tools.
\newblock In \textit{Advances in Neural Information Processing Systems (NeurIPS)}, 2023.

\bibitem[Shinn et al.(2023)]{shinn2023reflexion}
Shinn, N., Cassano, F., Labash, B., Gopinath, A., Narasimhan, K., and
Yao, S.
\newblock Reflexion: Language agents with verbal reinforcement learning.
\newblock In \textit{Advances in Neural Information Processing Systems (NeurIPS)}, 2023.

\bibitem[Shridhar et al.(2021)]{shridhar2021alfworld}
Shridhar, M., Yuan, X., C\^{o}t\'{e}, M.-A., Bisk, Y., Trischler, A., and
Hausknecht, M.
\newblock ALFWorld: Aligning text and embodied environments for interactive learning.
\newblock In \textit{International Conference on Learning Representations (ICLR)}, 2021.

\bibitem[Stone and Veloso(2000)]{stone2000}
Stone, P. and Veloso, M.
\newblock Multiagent systems: A survey from a machine learning perspective.
\newblock \textit{Autonomous Robots}, 8(3):345--383, 2000.

\bibitem[Yao et al.(2023)]{yao2023react}
Yao, S., Zhao, J., Yu, D., Du, N., Shafran, I., Narasimhan, K., and Cao, Y.
\newblock ReAct: Synergizing reasoning and acting in language models.
\newblock In \textit{International Conference on Learning Representations (ICLR)}, 2023.

\bibitem[Zhang et al.(2017)]{zhang2017ems}
Zhang, Z.-H. and Li, K.
\newblock A novel probabilistic formulation for locating and sizing emergency
medical service stations.
\newblock \textit{Annals of Operations Research}, 229(1):813--835, 2017.

\end{thebibliography}

\end{document}
